%! Sinusoidal Function Context and Data Modeling — Unit 6, Lesson 6.8 — Student Packet Cover Sheet
\documentclass[10pt]{article}
\usepackage{precalculus-article}
\usepackage{precalculus-boxes}
\usepackage{ltablex}
\keepXColumns

\begin{document}

% Full-bleed navy banner
\begin{tikzpicture}[remember picture, overlay]
  \fill[navy] (current page.north west)
    rectangle ([yshift=-0.9in]current page.north east);
\end{tikzpicture}
\vspace{-0.8in}

\begin{center}
  {\color{white}\textbf{\LARGE Precalculus}}\\[4pt]
  {\color{skymid}\large\textbf{Unit 6: Trigonometric Functions}}\\[6pt]
  {\color{white}\textbf{\large Lesson 6.8 \quad Sinusoidal Function Context and Data Modeling}}
\end{center}

\vspace{0.22in}
\namedateperiod
\vspace{0.18in}

\begin{learningtargetbox}
\small
\begin{enumerate}[leftmargin=1.5em, itemsep=5pt]
  \item I can determine the amplitude, vertical shift (midline), period, and phase shift of a
    sinusoidal model from a data table using the formulas
    $A = \tfrac{\max - \min}{2}$ and $D = \tfrac{\max + \min}{2}$. \hfill\textit{(LO 3.7.A)}
  \item I can write a sinusoidal function model $f(x) = A\sin(B(x-C))+D$ (or cosine form)
    that fits a set of periodic real-world data. \hfill\textit{(LO 3.7.A)}
  \item I can use a sinusoidal model to predict output values within the contextual domain
    and evaluate whether my predictions are reasonable. \hfill\textit{(LO 3.7.A)}
\end{enumerate}
\end{learningtargetbox}

\vspace{0.14in}

\begin{tocbox}
\small
\renewcommand{\arraystretch}{1.5}
\begin{tabularx}{\linewidth}{@{} c l X r @{}}
\rowcolor{sky}
\textbf{\#} & \textbf{Component} & \textbf{Description} & \textbf{Score} \\
\midrule
1 & Warm-Up        & Spiral review: reading sinusoidal graphs; computing amplitude and midline from data & \blank{1.2cm} \\
2 & Guided Notes   & Extracting parameters from data; writing and validating a sinusoidal model & \blank{1.2cm} \\
3 & Group Activity & Build and evaluate sinusoidal models from real-world data (tiered R/A/E) & \blank{1.2cm} \\
4 & Exit Ticket    & Identify amplitude, period, and midline from a small data table & \blank{1.2cm} \\
5 & Homework       & Mixed: parameter extraction, prediction, AP-style MC & \blank{1.2cm} \\
\midrule
  & \multicolumn{2}{r}{\textbf{Total}} & \blank{1.2cm} \\
\end{tabularx}
\end{tocbox}

\end{document}
